% GENERATED FILE. Edit canonical lessons, then run npm run generate:books.
% canonical-source-hash: fnv1a64:65833d424855a0c8
% canonical-lessons: ZH-C513-chuangye, ZH-C513-dadao, ZH-C513-zhengqu, ZH-C513-fendou, ZH-C513-qinshou, ZH-R513-review-509-513, ZH-R513-review-final-513

\chapter{Ambitions}
\label{ch:zh-ambitions}
\hlchaptermodality{513}

\begin{chapteropening}
\textbf{By the end of this chapter:} \emph{I can talk about starting a business, reaching a goal, striving for it and doing it with my own hands.}

Complete the last lesson of chapter 513: I can talk about starting a business, reaching a goal, striving for it and doing it with my own hands.
\end{chapteropening}

\section[chuàngyè]{chuàngyè — to start a business; to found an enterprise}

\label{lesson:ZH-C513-chuangye}



\begin{quote}
Before the new one: say the Mandarin for to become famous, then the Mandarin for a top school.
\end{quote}

\subsection*{You'll want to know: chuàngyè}
\emph{chuàngyè} — \textquotedblleft{}to start a business; to found an enterprise\textquotedblright{}.

\begin{quote}
\emph{dàxué bìyè yǐhòu wǒ xiǎng zìjǐ chuàngyè} — After university I want to start my own business.
\end{quote}

Its characters come later; for now it is a word for the ear.

\begin{cousinweb}
\emph{chuàngyè} is built from \emph{chuàng} 'to found' and \zh{业} 'an undertaking', both inherited from Middle Chinese, \emph{tsrhjangH} and \emph{ngjaep} (Guangyun, in Baxter's transcription). The dictionary gives the first as another spelling of a word it glosses, after an older dictionary, 'to found a craft and begin an undertaking'.
\end{cousinweb}

\begin{grammarlens}[title={What you've built}]
One more word for saying how it was.
\end{grammarlens}

\subsection*{Your turn}
\begin{itemize}
\raggedright
  \item \emph{Say it:} \emph{chuàngyè}
  \item \emph{Say it:} \emph{chuàngyè}, once more
  \item \emph{Say it:} \emph{dàxué bìyè yǐhòu wǒ xiǎng zìjǐ chuàngyè} — After university I want to start my own business.
  \item \emph{Recall:} say the Mandarin for to become famous, then the Mandarin for a top school, then say \emph{chuàngyè} again
\end{itemize}

\begin{tcolorbox}[breakable,colback=teal!4,colframe=teal!35!black,title={Before you move on}]
What does chuàngyè mean? (\textquotedblleft{}To start a business; to found an enterprise\textquotedblright{}.) Say it once more.
\end{tcolorbox}
\section[dádào]{dádào — to attain (a goal); to reach a level}

\label{lesson:ZH-C513-dadao}



\begin{quote}
Before the new one: say the Mandarin for a top school, then the Mandarin for to start a business.
\end{quote}

\subsection*{You'll want to know: dádào}
\emph{dádào} — \textquotedblleft{}to attain (a goal); to reach a level\textquotedblright{}.

\begin{quote}
\emph{tā xué le sān nián dádào le zìjǐ de mùbiāo} — He studied for three years and reached his goal.
\end{quote}

Its characters come later; for now it is a word for the ear.

\begin{cousinweb}
\emph{dádào} is built from \emph{dá} 'to get through' and \zh{到} 'to arrive', both inherited from Middle Chinese, \emph{dat} and \emph{tawH} (Guangyun, in Baxter's transcription), which glosses the first 'to reach through' and the second 'to arrive'.
\end{cousinweb}

\begin{grammarlens}[title={What you've built}]
One more word for saying how it was.
\end{grammarlens}

\subsection*{Your turn}
\begin{itemize}
\raggedright
  \item \emph{Say it:} \emph{dádào}
  \item \emph{Say it:} \emph{dádào}, once more
  \item \emph{Say it:} \emph{tā xué le sān nián dádào le zìjǐ de mùbiāo} — He studied for three years and reached his goal.
  \item \emph{Recall:} say the Mandarin for a top school, then the Mandarin for to start a business, then say \emph{dádào} again
\end{itemize}

\begin{tcolorbox}[breakable,colback=teal!4,colframe=teal!35!black,title={Before you move on}]
What does dádào mean? (\textquotedblleft{}To attain (a goal); to reach a level\textquotedblright{}.) Say it once more.
\end{tcolorbox}
\section[zhēngqǔ]{zhēngqǔ — to strive for; to fight for}

\label{lesson:ZH-C513-zhengqu}



\begin{quote}
Before the new one: say the Mandarin for to start a business, then the Mandarin for to attain (a goal).
\end{quote}

\subsection*{You'll want to know: zhēngqǔ}
\emph{zhēngqǔ} — \textquotedblleft{}to strive for; to fight for\textquotedblright{}.

\begin{quote}
\emph{wǒ yào zhēngqǔ kǎoshàng míngxiào} — I'm going to do my best to get into a top school.
\end{quote}

Its characters come later; for now it is a word for the ear.

\begin{cousinweb}
\emph{zhēngqǔ} is built from \emph{zhēng} 'to contend' and \zh{取} 'to take', both inherited from Middle Chinese, \emph{tsreang} and \emph{tshjuX} (Guangyun, in Baxter's transcription), which glosses the first 'to compete' and the second 'to receive'.
\end{cousinweb}

\begin{grammarlens}[title={What you've built}]
One more word for saying how it was.
\end{grammarlens}

\subsection*{Your turn}
\begin{itemize}
\raggedright
  \item \emph{Say it:} \emph{zhēngqǔ}
  \item \emph{Say it:} \emph{zhēngqǔ}, once more
  \item \emph{Say it:} \emph{wǒ yào zhēngqǔ kǎoshàng míngxiào} — I'm going to do my best to get into a top school.
  \item \emph{Recall:} say the Mandarin for to start a business, then the Mandarin for to attain (a goal), then say \emph{zhēngqǔ} again
\end{itemize}

\begin{tcolorbox}[breakable,colback=teal!4,colframe=teal!35!black,title={Before you move on}]
What does zhēngqǔ mean? (\textquotedblleft{}To strive for; to fight for\textquotedblright{}.) Say it once more.
\end{tcolorbox}
\section[fèndòu]{fèndòu — to struggle (towards a goal)}

\label{lesson:ZH-C513-fendou}



\begin{quote}
Before the new one: say the Mandarin for to attain (a goal), then the Mandarin for to strive for.
\end{quote}

\subsection*{You'll want to know: fèndòu}
\emph{fèndòu} — \textquotedblleft{}to struggle (towards a goal)\textquotedblright{}.

\begin{quote}
\emph{tā fèndòu le shí nián cái chénggōng} — He struggled for ten years before he succeeded.
\end{quote}

Its characters come later; for now it is a word for the ear.

\begin{cousinweb}
\emph{fèndòu} is built from \emph{fèn} 'to spread the wings' and \emph{dòu} 'to fight', both inherited from Middle Chinese, \emph{pjunH} and \emph{tuwH} (Guangyun, in Baxter's transcription), which glosses the first as a bird spreading its feathers and the second, after an older dictionary, as two men facing each other with weapons.
\end{cousinweb}

\begin{grammarlens}[title={What you've built}]
One more word for saying how it was.
\end{grammarlens}

\subsection*{Your turn}
\begin{itemize}
\raggedright
  \item \emph{Say it:} \emph{fèndòu}
  \item \emph{Say it:} \emph{fèndòu}, once more
  \item \emph{Say it:} \emph{tā fèndòu le shí nián cái chénggōng} — He struggled for ten years before he succeeded.
  \item \emph{Recall:} say the Mandarin for to attain (a goal), then the Mandarin for to strive for, then say \emph{fèndòu} again
\end{itemize}

\begin{tcolorbox}[breakable,colback=teal!4,colframe=teal!35!black,title={Before you move on}]
What does fèndòu mean? (\textquotedblleft{}To struggle (towards a goal)\textquotedblright{}.) Say it once more.
\end{tcolorbox}
\section[\texorpdfstring{\zh{亲手} (qīnshǒu)}{亲手 (qīnshǒu)}]{\zh{亲手} (qīnshǒu) — with one's own hands}

\label{lesson:ZH-C513-qinshou}



\begin{quote}
Before the new one: say the Mandarin for to strive for, then the Mandarin for to struggle (towards a goal).
\end{quote}

\subsection*{You'll want to know: \zh{亲手}}
\textbf{\zh{亲手}} — \emph{qīnshǒu} — \textquotedblleft{}with one's own hands\textquotedblright{}.

\begin{quote}
\textbf{\zh{我想亲手开一家饭店}} \emph{wǒ xiǎng qīnshǒu kāi yì jiā fàndiàn} — I want to open a restaurant with my own hands.
\end{quote}

\begin{cousinweb}
\textbf{\zh{亲手}} (\emph{qīnshǒu}) is built from \zh{亲} 'close' and \zh{手} 'a hand', both inherited from Middle Chinese, \emph{tshin} and \emph{syuwX} (Guangyun, in Baxter's transcription).
\end{cousinweb}

\begin{grammarlens}[title={What you've built}]
One more word for saying how it was.
\end{grammarlens}

\subsection*{Your turn}
\begin{itemize}
\raggedright
  \item \emph{Say it:} \emph{qīnshǒu}
  \item \emph{Say it:} \emph{qīnshǒu}, once more
  \item \emph{Say it:} \emph{wǒ xiǎng qīnshǒu kāi yì jiā fàndiàn} — I want to open a restaurant with my own hands.
  \item \emph{Recall:} say the Mandarin for to strive for, then the Mandarin for to struggle (towards a goal), then say \emph{qīnshǒu} again
\end{itemize}

\begin{tcolorbox}[breakable,colback=teal!4,colframe=teal!35!black,title={Before you move on}]
What does \zh{亲手} mean? (\textquotedblleft{}With one's own hands\textquotedblright{}.) Say it once more.
\end{tcolorbox}
\section[(dialogue)]{Looking back: from Parents to Ambitions}

\label{lesson:ZH-R513-review-509-513}



\begin{quote}
Say the Mandarin for to run the household and one's mother and father.
\end{quote}

\subsection*{Your turn}
\begin{itemize}
\raggedright
  \item \emph{Say it:} to run the household, one's mother and father, a mother and daughter, one's mother tongue, a stepmother
  \item \emph{Say it:} grandma on the mother's side, a grandson, a granddaughter, to dote on, to soothe a child
  \item \emph{Say it:} one's nearest and dearest, father (the formal word), mother (the formal word), one's own (by birth), personally (doing it oneself)
  \item \emph{Say it:} one's ambition in life, to pursue (a goal), to make something of oneself, to become famous, a top school
  \item \emph{Say it:} to start a business, to attain (a goal), to strive for, to struggle (towards a goal), with one's own hands
  \item \emph{Read it:} \zh{母} \zh{亲} — each one aloud, then a word that holds it
\end{itemize}

\begin{tcolorbox}[breakable,colback=teal!4,colframe=teal!35!black,title={Before you move on}]
To run the household? (\textbf{\zh{当家}}, \emph{dāngjiā}.) One's mother and father? (\textbf{\zh{父母}}, \emph{fùmǔ}.) A mother and daughter? (\textbf{\zh{母女}}, \emph{mǔnǚ}.) One's mother tongue? (\textbf{\zh{母语}}, \emph{mǔyǔ}.) A stepmother? (\textbf{\zh{后妈}}, \emph{hòumā}.) Grandma on the mother's side? (\emph{wàipó}.) A grandson? (\emph{sūnzi}.) A granddaughter? (\emph{sūnnǚ}.) To dote on? (\textbf{\zh{疼爱}}, \emph{téng'ài}.) To soothe a child? (\emph{hǒng}.) One's nearest and dearest? (\textbf{\zh{亲人}}, \emph{qīnrén}.) Father (the formal word)? (\textbf{\zh{父亲}}, \emph{fùqīn}.) Mother (the formal word)? (\textbf{\zh{母亲}}, \emph{mǔqīn}.) One's own (by birth)? (\textbf{\zh{亲生}}, \emph{qīnshēng}.) Personally (doing it oneself)? (\textbf{\zh{亲自}}, \emph{qīnzì}.) One's ambition in life? (\emph{zhìxiàng}.) To pursue (a goal)? (\emph{zhuīqiú}.) To make something of oneself? (\textbf{\zh{成才}}, \emph{chéngcái}.) To become famous? (\textbf{\zh{成名}}, \emph{chéngmíng}.) A top school? (\textbf{\zh{名校}}, \emph{míngxiào}.) To start a business? (\emph{chuàngyè}.) To attain (a goal)? (\emph{dádào}.) To strive for? (\emph{zhēngqǔ}.) To struggle (towards a goal)? (\emph{fèndòu}.) With one's own hands? (\textbf{\zh{亲手}}, \emph{qīnshǒu}.)
\end{tcolorbox}
\section[(dialogue)]{Ambitions — once more}

\label{lesson:ZH-R513-review-final-513}



\begin{quote}
Say the Mandarin for to start a business and to attain (a goal).
\end{quote}

\subsection*{Your turn}
\begin{itemize}
\raggedright
  \item \emph{Say it:} to start a business, to attain (a goal), to strive for, to struggle (towards a goal), with one's own hands
\end{itemize}

\begin{tcolorbox}[breakable,colback=teal!4,colframe=teal!35!black,title={Before you move on}]
To start a business? (\emph{chuàngyè}.) To attain (a goal)? (\emph{dádào}.) To strive for? (\emph{zhēngqǔ}.) To struggle (towards a goal)? (\emph{fèndòu}.) With one's own hands? (\textbf{\zh{亲手}}, \emph{qīnshǒu}.)
\end{tcolorbox}
